% !TEX root = ../main.tex
% --- LỜI CẢM ƠN ---
\newpage
\pagestyle{frontmatter}
\chapter*{LỜI CẢM ƠN}
\thispagestyle{frontmatter}

Khóa luận tốt nghiệp không chỉ là cột mốc khép lại chặng đường đại học mà còn là bước đệm quan trọng để chúng em vững tin bước vào môi trường làm việc thực tế. Để hoàn thành đề tài này, bên cạnh sự nỗ lực và tâm huyết của cả nhóm, chúng em đã nhận được sự động viên và hỗ trợ to lớn từ nhiều nguồn kiến thức quý báu.

Lời tri ân sâu sắc nhất, chúng em xin trân trọng gửi đến cô \textbf{ThS. Nguyễn Thị Hoàng Khánh}. Trong suốt thời gian thực hiện đề tài, Cô đã luôn tận tình định hướng, chỉ bảo và chia sẻ những kinh nghiệm thực tiễn vô cùng giá trị. Sự hướng dẫn sát sao của Cô chính là kim chỉ nam giúp chúng em khắc phục những thiếu sót và hoàn thiện sản phẩm một cách tốt nhất.

Chúng em cũng xin gửi lời cảm ơn chân thành đến Quý Thầy Cô trường Đại học Công nghiệp TP.HCM. Những kiến thức nền tảng và kỹ năng mềm mà Thầy Cô truyền đạt trong suốt 4 năm qua chính là hành trang vững chắc để chúng em thực hiện đồ án này cũng như phát triển sự nghiệp trong tương lai.

Cuối cùng, chúng em xin cảm ơn Ban Giám hiệu nhà trường đã luôn tạo điều kiện thuận lợi về cơ sở vật chất và môi trường học tập, giúp sinh viên phát huy tối đa năng lực của mình.

Chúng em xin chân thành cảm ơn!

\vspace{1cm}
\begin{flushright}
    \textit{TP. Hồ Chí Minh, ngày 20 tháng 4 năm 2026}\\
    \textbf{Sinh viên thực hiện}\\
    \textbf{Lai Thanh Sĩ}\\
    \textbf{Phạm Đức Tài}
\end{flushright}

% --- NHẬN XÉT GVHD ---
\newpage
\chapter*{NHẬN XÉT CỦA GIẢNG VIÊN HƯỚNG DẪN}
\thispagestyle{frontmatter}
\dotline \dotline \dotline \dotline \dotline
\dotline \dotline \dotline \dotline \dotline
\dotline \dotline \dotline 
\vspace{0.5cm}
\begin{flushright}
    \textit{TP. Hồ Chí Minh, ngày \dots tháng \dots năm 20\dots}\\
    \textbf{GIẢNG VIÊN HƯỚNG DẪN}\\
    \textit{(Ký và ghi rõ họ tên)}\\
    \vspace{1.5cm}
\end{flushright}

% --- NHẬN XÉT GVPB 1 ---
\newpage
\chapter*{NHẬN XÉT CỦA GIẢNG VIÊN PHẢN BIỆN 1}
\thispagestyle{frontmatter}
\dotline \dotline \dotline \dotline \dotline
\dotline \dotline \dotline \dotline \dotline
\dotline \dotline \dotline 
\vspace{0.5cm}
\begin{flushright}
    \textit{TP. Hồ Chí Minh, ngày \dots tháng \dots năm 20\dots}\\
    \textbf{GIẢNG VIÊN PHẢN BIỆN 1}\\
    \textit{(Ký và ghi rõ họ tên)}\\
    \vspace{1.5cm}
\end{flushright}

% --- NHẬN XÉT GVPB 2 ---
\newpage
\chapter*{NHẬN XÉT CỦA GIẢNG VIÊN PHẢN BIỆN 2}
\thispagestyle{frontmatter}
\dotline \dotline \dotline \dotline \dotline
\dotline \dotline \dotline \dotline \dotline
\dotline \dotline \dotline 
\vspace{0.5cm}
\begin{flushright}
    \textit{TP. Hồ Chí Minh, ngày \dots tháng \dots năm 20\dots}\\
    \textbf{GIẢNG VIÊN PHẢN BIỆN 2}\\
    \textit{(Ký và ghi rõ họ tên)}\\
    \vspace{1.5cm}
\end{flushright}

% --- MỤC LỤC ---
\clearpage
\pagestyle{frontmatter}
\begingroup
\let\clearpage\relax
\addcontentsline{toc}{chapter}{MỤC LỤC}
\tableofcontents
\endgroup
\thispagestyle{frontmatter}

% --- DANH MỤC HÌNH ẢNH ---
\clearpage
\begingroup
\let\clearpage\relax
\addcontentsline{toc}{chapter}{DANH MỤC CÁC HÌNH ẢNH}
\listoffigures
\endgroup
\thispagestyle{frontmatter}

% --- DANH MỤC BẢNG BIỂU ---
\clearpage
\begingroup
\let\clearpage\relax
\addcontentsline{toc}{chapter}{DANH MỤC CÁC BẢNG BIỂU}
\listoftables
\endgroup
\thispagestyle{frontmatter}

% --- DANH MỤC TỪ VIẾT TẮT ---
\clearpage
\chapter*{DANH MỤC CÁC TỪ VIẾT TẮT}
\addcontentsline{toc}{chapter}{DANH MỤC CÁC TỪ VIẾT TẮT}
\thispagestyle{frontmatter}

% Chia tỷ lệ các cột: STT(8) + Từ viết tắt(18) + Nghĩa đầy đủ(74) = 100\tbu
\setlength{\tbu}{\dimexpr(\linewidth - 6\tabcolsep - 4\arrayrulewidth)/100\relax}
\begin{longtable}{|>{\centering\arraybackslash}p{8\tbu}|p{18\tbu}|p{74\tbu}|}
    \hline
    \textbf{STT} & \textbf{Từ viết tắt} & \textbf{Nghĩa đầy đủ} \\
    \hline
    \endfirsthead
    \hline
    \textbf{STT} & \textbf{Từ viết tắt} & \textbf{Nghĩa đầy đủ} \\
    \hline
    \endhead
    1  & AES   & Automated Essay Scoring (Chấm điểm bài viết tự động) \\
    \hline
    2  & AI    & Artificial Intelligence (Trí tuệ nhân tạo) \\
    \hline
    3  & AMQP  & Advanced Message Queuing Protocol \\
    \hline
    4  & AMQPS & Advanced Message Queuing Protocol Secure (AMQP qua TLS) \\
    \hline
    5  & API   & Application Programming Interface \\
    \hline
    6  & APS   & Automated Pronunciation Scoring (Chấm điểm phát âm tự động) \\
    \hline
    7  & BaaS  & Backend as a Service \\
    \hline
    8  & CALL  & Computer-Assisted Language Learning (Học ngôn ngữ hỗ trợ máy tính) \\
    \hline
    9  & CDN   & Content Delivery Network (Mạng phân phối nội dung) \\
    \hline
    10 & CEFR  & Common European Framework of Reference for Languages \\
    \hline
    11 & CI/CD & Continuous Integration / Continuous Deployment (Tích hợp và triển khai liên tục) \\
    \hline
    12 & CNTT  & Công nghệ thông tin \\
    \hline
    13 & DI    & Dependency Injection (Tiêm phụ thuộc) \\
    \hline
    14 & DLQ   & Dead Letter Queue (Hàng đợi thư chết) \\
    \hline
    15 & DTO   & Data Transfer Object \\
    \hline
    16 & ERD   & Entity--Relationship Diagram (Sơ đồ thực thể -- quan hệ) \\
    \hline
    17 & FSRS  & Free Spaced Repetition Scheduler \\
    \hline
    18 & GCP   & Google Cloud Platform \\
    \hline
    19 & GCS   & Google Cloud Storage \\
    \hline
    20 & GVHD  & Giảng viên hướng dẫn \\
    \hline
    21 & GVPB  & Giảng viên phản biện \\
    \hline
    22 & HMAC  & Hash-based Message Authentication Code \\
    \hline
    23 & HTTP  & HyperText Transfer Protocol \\
    \hline
    24 & HTTPS & HyperText Transfer Protocol Secure \\
    \hline
    25 & IELTS & International English Language Testing System \\
    \hline
    26 & IPA   & International Phonetic Alphabet (Bảng phiên âm quốc tế) \\
    \hline
    27 & JSON  & JavaScript Object Notation \\
    \hline
    28 & JWT   & JSON Web Token \\
    \hline
    29 & LLM   & Large Language Model (Mô hình ngôn ngữ lớn) \\
    \hline
    30 & MAU   & Monthly Active Users (Người dùng hoạt động hàng tháng) \\
    \hline
    31 & NLP   & Natural Language Processing (Xử lý ngôn ngữ tự nhiên) \\
    \hline
    32 & ORM   & Object-Relational Mapping \\
    \hline
    33 & OTLP  & OpenTelemetry Protocol \\
    \hline
    34 & REST  & Representational State Transfer \\
    \hline
    35 & SDK   & Software Development Kit \\
    \hline
    36 & SQL   & Structured Query Language \\
    \hline
    37 & SRS   & Spaced Repetition System (Hệ thống lặp lại ngắt quãng) \\
    \hline
    38 & SSG   & Static Site Generation \\
    \hline
    39 & SSR   & Server-Side Rendering \\
    \hline
    40 & STT   & Speech-to-Text (Nhận dạng giọng nói thành văn bản) \\
    \hline
    41 & TLS   & Transport Layer Security \\
    \hline
    42 & TTL   & Time to Live \\
    \hline
    43 & UC    & Use Case (Ca sử dụng) \\
    \hline
    44 & UI    & User Interface (Giao diện người dùng) \\
    \hline
    45 & UX    & User Experience (Trải nghiệm người dùng) \\
    \hline
    46 & XP    & Experience Points (Điểm kinh nghiệm) \\
    \hline
    47 & ZPD   & Zone of Proximal Development (Vùng phát triển gần) \\
    \hline
\end{longtable}
